%==============================================================================
%  Hierarchical Clustering:
%  Background, Theory, Types and Worked Examples
%
%  Self-contained lecture notes for an introductory Machine Learning course
%  (B.Tech / M.Tech level).  Prerequisites: basic linear algebra and distances.
%
%  Compile with:  pdflatex hierarchical_clustering.tex   (run twice for refs)
%  Requires a full TeX distribution (TeX Live / MiKTeX) with pgfplots + tikz.
%==============================================================================
\documentclass[11pt,a4paper]{report}

%------------------------------- Page geometry --------------------------------
\usepackage[a4paper,margin=1in]{geometry}

%------------------------------- Mathematics ----------------------------------
\usepackage{amsmath,amssymb,amsthm,mathtools}
\usepackage{bm}

%------------------------------- Tables & layout ------------------------------
\usepackage{booktabs}
\usepackage{array}
\usepackage{multirow}
\usepackage{enumitem}
\usepackage{caption}

%------------------------------- Graphics -------------------------------------
\usepackage{graphicx}
\usepackage{xcolor}
\usepackage{tikz}
\usepackage{pgfplots}
\pgfplotsset{compat=1.18}
\usetikzlibrary{arrows.meta,calc}

%------------------------------- Code listings --------------------------------
\usepackage{listings}

%------------------------------- Hyperlinks (load last) -----------------------
\usepackage[colorlinks=true,linkcolor=blue!60!black,citecolor=blue!60!black,
            urlcolor=blue!60!black,linktocpage=true]{hyperref}

%------------------------------- Theorem-like envs ----------------------------
\theoremstyle{definition}
\newtheorem{definition}{Definition}[chapter]
\newtheorem{example}{Example}[chapter]
\newtheorem{exercise}{Exercise}[chapter]
\theoremstyle{plain}
\newtheorem{theorem}{Theorem}[chapter]
\theoremstyle{remark}
\newtheorem*{remark}{Remark}

%------------------------------- Handy macros ---------------------------------
\newcommand{\R}{\mathbb{R}}
\newcommand{\vx}{\mathbf{x}}
\newcommand{\vy}{\mathbf{y}}
\newcommand{\va}{\mathbf{a}}
\newcommand{\vb}{\mathbf{b}}
\newcommand{\C}{\mathcal{C}}
\DeclareMathOperator*{\argmin}{arg\,min}
\DeclareMathOperator*{\argmax}{arg\,max}

%------------------------------- Listing style --------------------------------
\definecolor{codegreen}{rgb}{0.0,0.5,0.0}
\definecolor{codegray}{rgb}{0.5,0.5,0.5}
\definecolor{codeblue}{rgb}{0.1,0.1,0.6}
\definecolor{backcol}{rgb}{0.97,0.97,0.95}
\lstdefinestyle{python}{
  backgroundcolor=\color{backcol},language=Python,
  commentstyle=\color{codegreen},keywordstyle=\color{codeblue}\bfseries,
  numberstyle=\tiny\color{codegray},stringstyle=\color{codegreen},
  basicstyle=\ttfamily\footnotesize,breaklines=true,captionpos=b,
  numbers=left,numbersep=6pt,showstringspaces=false,frame=single,
  rulecolor=\color{codegray},tabsize=2}

%------------------------------- Dendrogram helper ----------------------------
% \dnode draws a horizontal merge bar joining two vertical stems.
% arguments: (x1,y1)(x2,y2) merge-height  ->  merge at given height h
\newcommand{\merge}[5]{% x1 y1 x2 y2 h
  \draw[thick] (#1,#2) -- (#1,#5) -- (#3,#5) -- (#3,#4);
}

%==============================================================================
\begin{document}

%------------------------------- Title ----------------------------------------
\begin{titlepage}
\centering
\vspace*{2cm}
{\Huge\bfseries Hierarchical Clustering\par}
\vspace{0.6cm}
{\Large Background, Theory, Types\\[2pt]
and Worked Examples by Hand\par}
\vspace{1.5cm}
{\large A Self-Contained Set of Lecture Notes\par}
\vspace{0.4cm}
{\large Introduction to Machine Learning\par}
\vspace{2.5cm}
\begin{tabular}{c}
\hline\\[-6pt]
\textbf{Prerequisites:} Basic Linear Algebra and Distance Measures\\[4pt]
\textbf{Level:} B.Tech / M.Tech\\[4pt]
\hline
\end{tabular}
\vfill
{\large \today\par}
\end{titlepage}

\tableofcontents
\clearpage

%==============================================================================
\chapter{Background: Clustering in Machine Learning}
%==============================================================================

\section{Supervised vs.\ Unsupervised Learning}
Recall that in \emph{supervised} learning every training example carries a
label, and the goal is to learn a mapping from inputs to those labels
(e.g.\ linear regression, classification). In \emph{unsupervised} learning we
are given only the inputs
\[
\{\vx^{(1)},\vx^{(2)},\dots,\vx^{(m)}\},\qquad \vx^{(i)}\in\R^{n},
\]
with \textbf{no target values}. The aim is to discover hidden structure in the
data. \emph{Clustering} is the most important unsupervised task.

\section{What is Clustering?}
\begin{definition}[Clustering]
Clustering is the task of partitioning a set of objects into groups
(\emph{clusters}) such that objects in the same cluster are \emph{similar} to
one another and objects in different clusters are \emph{dissimilar}.
\end{definition}
Similarity is quantified by a \emph{distance} (or dissimilarity) measure: small
distance $\Rightarrow$ high similarity. A good clustering therefore minimizes
intra-cluster distance and maximizes inter-cluster distance.

\section{Applications of Clustering}
\begin{itemize}[nosep]
  \item \textbf{Customer segmentation} --- grouping shoppers by behaviour.
  \item \textbf{Document / topic grouping} --- organizing news or papers.
  \item \textbf{Bioinformatics} --- grouping genes with similar expression;
        building phylogenetic (evolutionary) trees.
  \item \textbf{Image segmentation} --- grouping pixels into regions.
  \item \textbf{Anomaly detection} --- points that join no cluster are outliers.
  \item \textbf{Taxonomies} --- organizing species, products, or files into a
        nested hierarchy.
\end{itemize}

\section{Partitional vs.\ Hierarchical Clustering}
Clustering algorithms fall into two broad families.
\begin{center}
\begin{tabular}{@{}p{6.4cm}p{6.4cm}@{}}
\toprule
\textbf{Partitional (flat)} & \textbf{Hierarchical}\\
\midrule
Produces a single partition into $k$ clusters. & Produces a nested tree of
clusters at all levels.\\
$k$ must be chosen in advance. & $k$ need not be fixed beforehand.\\
Example: $k$-means, $k$-medoids. & Example: agglomerative, divisive.\\
Output: cluster assignments. & Output: a \emph{dendrogram}.\\
\bottomrule
\end{tabular}
\end{center}

\noindent
\textbf{Hierarchical clustering} is the subject of these notes. Its great
advantages are that (i) it needs no pre-specified number of clusters, and
(ii) it yields a rich tree (dendrogram) that can be ``cut'' at any level to
obtain any desired number of clusters.

%==============================================================================
\chapter{Distance and Similarity Measures}
%==============================================================================

Every hierarchical method begins from pairwise distances between the objects we
wish to cluster. We represent each object as a single \emph{point} in
$n$-dimensional space, where $n$ is the number of features (measurements)
describing it.

To state a distance formula we need names for \emph{two arbitrary points}. We
deliberately do \textbf{not} call them $x$ and $y$: in the plane those letters
already denote the two \emph{coordinates} of one point, so reusing them as the
\emph{names} of two points is exactly what causes confusion. Instead we name the
two points $\va$ and $\vb$---just as the worked examples name the objects
$A,B,C,\dots$. Being points of $\R^n$, each one carries $n$ coordinates (one per
feature), which we write with a subscript on the point's own letter:
\[
\va=(a_1,a_2,\dots,a_n),\qquad \vb=(b_1,b_2,\dots,b_n).
\]
Here $a_k$ is the $k$-th coordinate of the \emph{single} point $\va$; it is
\textbf{not} a separate point. In the familiar plane ($n=2$) this is just the
usual abscissa/ordinate pair: $\va=(a_1,a_2)$ is one point and
$\vb=(b_1,b_2)$ is another, and the formulas below reduce to the schoolbook
$\sqrt{(a_1-b_1)^2+(a_2-b_2)^2}$. For $n>2$ we simply have more coordinates per
point.

\section{Common Distance Measures}
\begin{definition}[Euclidean distance ($L_2$)]
\begin{equation}
d(\va,\vb)=\sqrt{\sum_{k=1}^{n}(a_k-b_k)^2}.
\label{eq:euclid}
\end{equation}
\end{definition}
\begin{definition}[Manhattan distance ($L_1$, city-block)]
\[
d(\va,\vb)=\sum_{k=1}^{n}\lvert a_k-b_k\rvert .
\]
\end{definition}
\begin{definition}[Minkowski distance ($L_p$)]
\[
d(\va,\vb)=\Bigl(\sum_{k=1}^{n}\lvert a_k-b_k\rvert^{p}\Bigr)^{1/p},
\]
which gives Manhattan for $p=1$ and Euclidean for $p=2$.
\end{definition}
\begin{definition}[Cosine dissimilarity]
\[
d(\va,\vb)=1-\frac{\va^{\!\top}\vb}{\lVert\va\rVert\,\lVert\vb\rVert},
\]
useful for text/high-dimensional data where direction matters more than
magnitude.
\end{definition}

A valid distance (a \emph{metric}) satisfies non-negativity, symmetry
$d(\va,\vb)=d(\vb,\va)$, identity $d(\va,\va)=0$, and the triangle inequality.

\section{The Proximity (Distance) Matrix}
For $m$ objects the pairwise distances are stored in a symmetric
$m\times m$ \emph{proximity matrix} $\mathbf{D}$ with $D_{ij}=d(\vx^{(i)},\vx^{(j)})$
and zero diagonal:
\[
\mathbf{D}=\begin{bmatrix}
0 & D_{12} & \cdots & D_{1m}\\
D_{12} & 0 & \cdots & D_{2m}\\
\vdots & & \ddots & \vdots\\
D_{1m} & D_{2m} & \cdots & 0
\end{bmatrix}.
\]
Because $\mathbf D$ is symmetric, we only ever need its lower (or upper)
triangle. \textbf{This matrix is the sole input to agglomerative hierarchical
clustering}---the raw coordinates are not needed once $\mathbf D$ is formed.

\begin{remark}[Feature scaling]
Distances are dominated by large-range features. Standardizing each feature to
zero mean and unit variance before computing $\mathbf D$ prevents any single
feature from dominating.
\end{remark}

%==============================================================================
\chapter{Hierarchical Clustering: Theory}
\label{ch:theory}
%==============================================================================

\section{Core Idea}
Hierarchical clustering builds a \emph{hierarchy} of clusters, represented as a
binary tree called a \textbf{dendrogram}. Each leaf is a single object; each
internal node is a cluster formed by merging its two children; the height of a
node is the distance at which the merge happened.

\section{The Dendrogram}
A dendrogram encodes the entire clustering at every granularity simultaneously.
Cutting it with a horizontal line at height $h$ severs all links above $h$ and
leaves a specific number of clusters: a low cut yields many small clusters, a
high cut yields few large clusters (Figure~\ref{fig:cut}).

\begin{figure}[ht]
\centering
\begin{tikzpicture}[yscale=0.55,xscale=1.0]
% simple illustrative dendrogram
\foreach \x/\l in {1/a,2/b,3/c,4/d,5/e} \node[below] at (\x,0){\l};
\merge{1}{0}{2}{0}{2}          % a,b at 2
\merge{4}{0}{5}{0}{3}          % d,e at 3
\merge{1.5}{2}{3}{0}{4}        % (ab),c at 4
\merge{2.25}{4}{4.5}{3}{6}     % (abc),(de) at 6
% cut line
\draw[red,dashed,thick] (0.3,5) -- (5.7,5) node[right]{\small cut at $h{=}5$};
\draw[->] (-0.4,0)--(-0.4,7) node[above]{\scriptsize height};
\foreach \y in {2,3,4,6} \draw (-0.5,\y)--(-0.3,\y) node[left]{\scriptsize \y};
\end{tikzpicture}
\caption{A dendrogram. The dashed red cut at $h=5$ produces two clusters:
$\{a,b,c\}$ and $\{d,e\}$.}
\label{fig:cut}
\end{figure}

\section{The Generic Agglomerative Algorithm}
The \emph{agglomerative} (bottom-up) procedure is by far the most common.

\begin{center}
\fbox{\parbox{0.9\textwidth}{
\textbf{Agglomerative Hierarchical Clustering}
\begin{enumerate}[nosep,leftmargin=2.2em]
  \item Start with each object as its own cluster ($m$ clusters).
  \item Compute the proximity matrix $\mathbf D$ between all clusters.
  \item Find the pair of clusters with the \emph{smallest} distance and
        \textbf{merge} them into one.
  \item \textbf{Update} $\mathbf D$: recompute distances from the new cluster
        to all others using a chosen \emph{linkage} rule.
  \item Repeat steps 3--4 until a single cluster remains.
  \item Record every merge and its height to draw the dendrogram.
\end{enumerate}
}}
\end{center}

\noindent
Steps 3--4 are repeated exactly $m-1$ times. The only design choice is the
\emph{linkage rule} in step 4, which defines the distance between two
\emph{clusters} (not just two points). Different linkages give different trees;
they are the ``types'' studied next.

%==============================================================================
\chapter{Types of Hierarchical Clustering}
\label{ch:types}
%==============================================================================

\section{Agglomerative vs.\ Divisive}
\begin{center}
\begin{tabular}{@{}p{6.4cm}p{6.4cm}@{}}
\toprule
\textbf{Agglomerative (AGNES, bottom-up)} & \textbf{Divisive (DIANA, top-down)}\\
\midrule
Start: every object is its own cluster. & Start: all objects in one cluster.\\
Repeatedly \emph{merge} the two closest clusters. & Repeatedly \emph{split} the
most heterogeneous cluster.\\
$m-1$ merges. & Up to $m-1$ splits.\\
Cheaper and far more common. & Expensive ($2^{m-1}$ ways to split a cluster);
rarely used.\\
\bottomrule
\end{tabular}
\end{center}
The remaining chapters focus on the agglomerative family, whose ``types'' are
distinguished by the linkage criterion.

\section{Linkage Criteria}
Let $\C_i$ and $\C_j$ be two clusters. The distance
$d(\C_i,\C_j)$ between them can be defined in several ways.

\begin{description}[leftmargin=0pt,itemsep=4pt]
\item[\textbf{Single linkage (MIN / nearest neighbour).}]
The distance is that of the \emph{closest} pair of points, one from each
cluster:
\begin{equation}
d_{\text{single}}(\C_i,\C_j)=\min_{\va\in\C_i,\ \vb\in\C_j} d(\va,\vb).
\label{eq:single}
\end{equation}
Tends to produce long, ``chained'' clusters; sensitive to noise bridges.

\item[\textbf{Complete linkage (MAX / farthest neighbour).}]
The distance is that of the \emph{farthest} pair:
\begin{equation}
d_{\text{complete}}(\C_i,\C_j)=\max_{\va\in\C_i,\ \vb\in\C_j} d(\va,\vb).
\label{eq:complete}
\end{equation}
Produces compact, roughly spherical clusters; resists chaining.

\item[\textbf{Average linkage (group average, UPGMA).}]
The mean of all inter-cluster pairwise distances:
\begin{equation}
d_{\text{average}}(\C_i,\C_j)=\frac{1}{|\C_i|\,|\C_j|}
   \sum_{\va\in\C_i}\sum_{\vb\in\C_j} d(\va,\vb).
\label{eq:average}
\end{equation}
A compromise between single and complete linkage.

\item[\textbf{Centroid linkage.}]
The distance between the cluster centroids
$\bm\mu_i,\bm\mu_j$: $\;d(\C_i,\C_j)=\lVert\bm\mu_i-\bm\mu_j\rVert$.

\item[\textbf{Ward's method (minimum variance).}]
Merge the pair whose union causes the smallest increase in total
within-cluster sum of squares (variance). It tends to produce clusters of
similar size and is very popular in practice.
\end{description}

\section{The Lance--Williams Update Formula}
All of the above linkages can be updated \emph{without} revisiting the original
points, using one recurrence. When clusters $\C_i$ and $\C_j$ merge into
$\C_{ij}$, the distance to any other cluster $\C_k$ is
\begin{equation}
d(\C_{ij},\C_k)=\alpha_i\,d(\C_i,\C_k)+\alpha_j\,d(\C_j,\C_k)
              +\beta\,d(\C_i,\C_j)+\gamma\,\lvert d(\C_i,\C_k)-d(\C_j,\C_k)\rvert .
\label{eq:lw}
\end{equation}
Choosing the coefficients reproduces each linkage (with
$n_i=|\C_i|$ etc.):
\begin{center}
\begin{tabular}{@{}lcccc@{}}
\toprule
Linkage & $\alpha_i$ & $\alpha_j$ & $\beta$ & $\gamma$\\
\midrule
Single   & $\tfrac12$ & $\tfrac12$ & $0$ & $-\tfrac12$\\[2pt]
Complete & $\tfrac12$ & $\tfrac12$ & $0$ & $+\tfrac12$\\[2pt]
Average  & $\dfrac{n_i}{n_i+n_j}$ & $\dfrac{n_j}{n_i+n_j}$ & $0$ & $0$\\[6pt]
Ward     & $\dfrac{n_i+n_k}{n_i+n_j+n_k}$ & $\dfrac{n_j+n_k}{n_i+n_j+n_k}$
         & $\dfrac{-n_k}{n_i+n_j+n_k}$ & $0$\\
\bottomrule
\end{tabular}
\end{center}
For example, single linkage sets
$d(\C_{ij},\C_k)=\tfrac12 d_{ik}+\tfrac12 d_{jk}-\tfrac12|d_{ik}-d_{jk}|
=\min(d_{ik},d_{jk})$, recovering \eqref{eq:single}. This is exactly the
``take the minimum'' shortcut we will use by hand.

%==============================================================================
\chapter{Worked Example 1: Single Linkage (by hand)}
\label{ch:single}
%==============================================================================

\section{The Data}
We are given five objects $A,B,C,D,E$ and the following symmetric proximity
matrix (only the lower triangle is needed). These could be Euclidean distances
between five cities, five gene profiles, etc.
\begin{equation}
\mathbf{D}^{(0)}=
\begin{array}{c|ccccc}
   & A & B & C & D & E\\ \hline
A  & 0 & 9 & 3 & 6 & 11\\
B  & 9 & 0 & 7 & 5 & 10\\
C  & 3 & 7 & 0 & 9 & 2\\
D  & 6 & 5 & 9 & 0 & 8\\
E  & 11& 10& 2 & 8 & 0
\end{array}
\label{eq:D0}
\end{equation}
Under \textbf{single linkage} the distance between two clusters is the
\emph{minimum} of the pairwise distances, so at each merge we replace the two
merged rows/columns by their element-wise minimum.

\section{Step 1 --- First Merge}
Scan the matrix for the smallest off-diagonal entry. The minimum is
$D_{CE}=2$. \textbf{Merge $C$ and $E$} into cluster $(CE)$ at height $2$.
Update using $d\bigl((CE),X\bigr)=\min\{d(C,X),d(E,X)\}$:
\[
\begin{aligned}
d((CE),A)&=\min(3,11)=3, &\quad d((CE),B)&=\min(7,10)=7,\\
d((CE),D)&=\min(9,8)=8. &&
\end{aligned}
\]
\begin{equation}
\mathbf{D}^{(1)}=
\begin{array}{c|cccc}
    & A & B & D & (CE)\\ \hline
A   & 0 & 9 & 6 & 3\\
B   & 9 & 0 & 5 & 7\\
D   & 6 & 5 & 0 & 8\\
(CE)& 3 & 7 & 8 & 0
\end{array}
\end{equation}

\section{Step 2 --- Second Merge}
Smallest entry in $\mathbf D^{(1)}$ is $d(A,(CE))=3$. \textbf{Merge $A$ with
$(CE)$} into $(ACE)$ at height $3$. Update by minima:
\[
d((ACE),B)=\min(9,7)=7,\qquad d((ACE),D)=\min(6,8)=6.
\]
\begin{equation}
\mathbf{D}^{(2)}=
\begin{array}{c|ccc}
     & B & D & (ACE)\\ \hline
B    & 0 & 5 & 7\\
D    & 5 & 0 & 6\\
(ACE)& 7 & 6 & 0
\end{array}
\end{equation}

\section{Step 3 --- Third Merge}
Smallest entry is $d(B,D)=5$. \textbf{Merge $B$ and $D$} into $(BD)$ at height
$5$. Update:
\[
d((BD),(ACE))=\min\bigl(d(B,(ACE)),d(D,(ACE))\bigr)=\min(7,6)=6.
\]
\begin{equation}
\mathbf{D}^{(3)}=
\begin{array}{c|cc}
     & (ACE) & (BD)\\ \hline
(ACE)& 0 & 6\\
(BD) & 6 & 0
\end{array}
\end{equation}

\section{Step 4 --- Final Merge}
Only two clusters remain; merge $(ACE)$ and $(BD)$ at height $6$ into the root
$\{A,B,C,D,E\}$.

\section{Merge Summary and Dendrogram}
\begin{center}
\begin{tabular}{@{}clcc@{}}
\toprule
Step & Clusters merged & Height & \# clusters after\\
\midrule
1 & $C,\;E$          & $2$ & 4\\
2 & $A,\;(CE)$       & $3$ & 3\\
3 & $B,\;D$          & $5$ & 2\\
4 & $(ACE),\;(BD)$   & $6$ & 1\\
\bottomrule
\end{tabular}
\end{center}

\begin{figure}[ht]
\centering
\begin{tikzpicture}[yscale=0.65,xscale=1.15]
\foreach \x/\l in {1/A,2/C,3/E,4/B,5/D} \node[below] at (\x,0){\l};
\merge{2}{0}{3}{0}{2}         % C,E at 2  -> center x=2.5
\merge{1}{0}{2.5}{2}{3}       % A,(CE) at 3 -> center x=1.75
\merge{4}{0}{5}{0}{5}         % B,D at 5  -> center x=4.5
\merge{1.75}{3}{4.5}{5}{6}    % (ACE),(BD) at 6
\draw[->] (0.2,0)--(0.2,7.2) node[above]{\scriptsize height};
\foreach \y in {2,3,5,6} \draw (0.1,\y)--(0.3,\y) node[left]{\scriptsize \y};
\end{tikzpicture}
\caption{Single-linkage dendrogram. Note the ``chaining'': $A$ attaches to the
$(CE)$ group early, then the two big groups join last.}
\label{fig:singdend}
\end{figure}

\begin{remark}[Reading a cut]
Cutting Figure~\ref{fig:singdend} just below height $6$ yields two clusters
$\{A,C,E\}$ and $\{B,D\}$; cutting below $5$ yields three: $\{A,C,E\}$,
$\{B\}$, $\{D\}$.
\end{remark}

%==============================================================================
\chapter{Worked Example 2: Complete Linkage (by hand)}
\label{ch:complete}
%==============================================================================

We repeat the procedure on the \emph{same} data \eqref{eq:D0}, but now with
\textbf{complete linkage}, where cluster distance is the \emph{maximum} of the
pairwise distances. The only change is that updates take the element-wise
\emph{maximum}.

\section{Step 1}
Minimum entry is still $D_{CE}=2$: \textbf{merge $C,E$} at height $2$. Update by
\emph{maxima}, $d((CE),X)=\max\{d(C,X),d(E,X)\}$:
\[
d((CE),A)=\max(3,11)=11,\quad d((CE),B)=\max(7,10)=10,\quad
d((CE),D)=\max(9,8)=9.
\]
\begin{equation}
\mathbf{D}^{(1)}=
\begin{array}{c|cccc}
    & A & B & D & (CE)\\ \hline
A   & 0 & 9 & 6 & 11\\
B   & 9 & 0 & 5 & 10\\
D   & 6 & 5 & 0 & 9\\
(CE)&11 &10 & 9 & 0
\end{array}
\end{equation}

\section{Step 2}
Smallest remaining entry is $d(B,D)=5$: \textbf{merge $B,D$} at height $5$.
Update by maxima:
\[
d((BD),A)=\max(9,6)=9,\qquad d((BD),(CE))=\max(10,9)=10.
\]
\begin{equation}
\mathbf{D}^{(2)}=
\begin{array}{c|ccc}
     & A & (CE) & (BD)\\ \hline
A    & 0 & 11 & 9\\
(CE) &11 & 0  & 10\\
(BD) & 9 & 10 & 0
\end{array}
\end{equation}
Notice that complete linkage merges $B,D$ \emph{second} here, whereas single
linkage merged $A$ with $(CE)$ second---the trees already diverge.

\section{Step 3}
Smallest entry is $d(A,(BD))=9$: \textbf{merge $A$ with $(BD)$} into $(ABD)$ at
height $9$. Update:
\[
d((ABD),(CE))=\max\bigl(d(A,(CE)),d((BD),(CE))\bigr)=\max(11,10)=11.
\]
\begin{equation}
\mathbf{D}^{(3)}=
\begin{array}{c|cc}
     & (ABD) & (CE)\\ \hline
(ABD)& 0 & 11\\
(CE) &11 & 0
\end{array}
\end{equation}

\section{Step 4}
Merge the last two clusters $(ABD)$ and $(CE)$ at height $11$.

\section{Merge Summary and Dendrogram}
\begin{center}
\begin{tabular}{@{}clcc@{}}
\toprule
Step & Clusters merged & Height & \# clusters after\\
\midrule
1 & $C,\;E$          & $2$  & 4\\
2 & $B,\;D$          & $5$  & 3\\
3 & $A,\;(BD)$       & $9$  & 2\\
4 & $(ABD),\;(CE)$   & $11$ & 1\\
\bottomrule
\end{tabular}
\end{center}

\begin{figure}[ht]
\centering
\begin{tikzpicture}[yscale=0.38,xscale=1.15]
\foreach \x/\l in {1/C,2/E,3/A,4/B,5/D} \node[below] at (\x,0){\l};
\merge{1}{0}{2}{0}{2}         % C,E at 2 -> x=1.5
\merge{4}{0}{5}{0}{5}         % B,D at 5 -> x=4.5
\merge{3}{0}{4.5}{5}{9}       % A,(BD) at 9 -> x=3.75
\merge{1.5}{2}{3.75}{9}{11}   % (CE),(ABD) at 11
\draw[->] (0.2,0)--(0.2,12.5) node[above]{\scriptsize height};
\foreach \y in {2,5,9,11} \draw (0.1,\y)--(0.3,\y) node[left]{\scriptsize \y};
\end{tikzpicture}
\caption{Complete-linkage dendrogram on the same data. The merge heights are
larger and the tree is more balanced/compact than the single-linkage tree.}
\label{fig:compdend}
\end{figure}

%==============================================================================
\chapter{Worked Example 3: Average Linkage (by hand)}
\label{ch:average}
%==============================================================================

Average linkage uses the \emph{mean} of all pairwise inter-cluster distances,
equation \eqref{eq:average}. When a cluster of size $|\C_i|$ meets a cluster of
size $|\C_j|$ there are $|\C_i|\,|\C_j|$ pairs to average.

\section{Steps}
\textbf{Step 1.} Minimum is $D_{CE}=2$: merge $C,E$ (size 2) at height $2$.
Average of the two singleton distances:
\[
d((CE),A)=\tfrac{3+11}{2}=7,\quad
d((CE),B)=\tfrac{7+10}{2}=8.5,\quad
d((CE),D)=\tfrac{9+8}{2}=8.5.
\]
\begin{equation}
\mathbf{D}^{(1)}=
\begin{array}{c|cccc}
    & A & B & D & (CE)\\ \hline
A   & 0 & 9 & 6 & 7\\
B   & 9 & 0 & 5 & 8.5\\
D   & 6 & 5 & 0 & 8.5\\
(CE)& 7 &8.5&8.5& 0
\end{array}
\end{equation}

\textbf{Step 2.} Minimum is $d(B,D)=5$: merge $B,D$ at height $5$.
\[
d((BD),A)=\tfrac{9+6}{2}=7.5,\qquad
d((BD),(CE))=\frac{d_{BC}+d_{BE}+d_{DC}+d_{DE}}{4}
=\frac{7+10+9+8}{4}=8.5 .
\]
\begin{equation}
\mathbf{D}^{(2)}=
\begin{array}{c|ccc}
     & A & (CE) & (BD)\\ \hline
A    & 0 & 7   & 7.5\\
(CE) & 7 & 0   & 8.5\\
(BD) &7.5& 8.5 & 0
\end{array}
\end{equation}

\textbf{Step 3.} Minimum is $d(A,(CE))=7$: merge $A$ with $(CE)$ into $(ACE)$ at
height $7$. Averaging over the $3\times2=6$ pairs between $\{A,C,E\}$ and
$\{B,D\}$:
\[
d((ACE),(BD))=\frac{d_{AB}+d_{AD}+d_{CB}+d_{CD}+d_{EB}+d_{ED}}{6}
=\frac{9+6+7+9+10+8}{6}=\frac{49}{6}\approx 8.17 .
\]
\textbf{Step 4.} Merge $(ACE)$ and $(BD)$ at height $\approx 8.17$.

\section{Merge Summary}
\begin{center}
\begin{tabular}{@{}clc@{}}
\toprule
Step & Clusters merged & Height\\
\midrule
1 & $C,\;E$        & $2$\\
2 & $B,\;D$        & $5$\\
3 & $A,\;(CE)$     & $7$\\
4 & $(ACE),\;(BD)$ & $49/6\approx8.17$\\
\bottomrule
\end{tabular}
\end{center}
The order of merges here matches complete linkage in its first two steps but
matches single linkage in step~3 ($A$ joins $(CE)$)---illustrating how average
linkage sits ``between'' the two extremes.

%==============================================================================
\chapter{Worked Example 4: Divisive Clustering --- DIANA (by hand)}
\label{ch:diana}
%==============================================================================

All previous examples were \emph{agglomerative} (bottom-up). We now work the
\textbf{divisive} (top-down) approach on the \emph{same} data \eqref{eq:D0}
using \textbf{DIANA} (DIvisive ANAlysis, Kaufman \& Rousseeuw). Divisive
clustering starts with every object in one big cluster and repeatedly
\emph{splits} a cluster into two.

\section{The DIANA Algorithm}
The subtlety of divisive clustering is that a cluster of $p$ objects can be cut
into two non-empty parts in $2^{p-1}-1$ ways---far too many to test. DIANA
avoids this with a clever \emph{splinter-group} heuristic.

\begin{center}
\fbox{\parbox{0.92\textwidth}{
\textbf{DIANA (one split of a cluster $\C$)}
\begin{enumerate}[nosep,leftmargin=2.2em]
  \item \textbf{Seed.} Find the object with the largest average dissimilarity
        to the rest of $\C$; move it out to start a new \emph{splinter group}
        $\C_{\text{new}}$. The rest form the \emph{old party} $\C_{\text{old}}$.
  \item \textbf{Reassign.} For every object $x\in\C_{\text{old}}$ compute
        \[
        D_x=\underbrace{a(x,\C_{\text{old}}\setminus\{x\})}_{\text{avg dist to old party}}
            -\underbrace{a(x,\C_{\text{new}})}_{\text{avg dist to splinter}} ,
        \]
        where $a(x,S)=\frac{1}{|S|}\sum_{y\in S}d(x,y)$ is the average distance
        from $x$ to set $S$.
  \item If $\max_x D_x>0$, move that object to the splinter group and repeat
        step~2. Otherwise \textbf{stop}: the split is
        $\C\to\C_{\text{old}}\mid\C_{\text{new}}$.
\end{enumerate}
\textbf{Which cluster to split next?} Among all current clusters, pick the one
with the largest \emph{diameter} (its maximum internal pairwise distance).
Repeat until every object is a singleton (or a stopping level is reached).
}}
\end{center}

\noindent
Intuitively, $D_x>0$ means $x$ is on average \emph{closer} to the breakaway
splinter group than to the old party, so it should defect.

\section{Split 1 --- Divide the Root $\{A,B,C,D,E\}$}
The whole set has diameter $\max_{ij}D_{ij}=11$ (the pair $A,E$), so it is the
cluster to split.

\subsection*{Seeding the splinter group}
Compute each object's average dissimilarity to the other four (row averages of
\eqref{eq:D0}):
\begin{center}
\begin{tabular}{@{}cccccc@{}}
\toprule
Object & $A$ & $B$ & $C$ & $D$ & $E$\\
\midrule
Avg.\ dissimilarity & $\tfrac{29}{4}{=}7.25$ & $\tfrac{31}{4}{=}7.75$
 & $\tfrac{21}{4}{=}5.25$ & $\tfrac{28}{4}{=}7.00$ & $\tfrac{31}{4}{=}7.75$\\
\bottomrule
\end{tabular}
\end{center}
The maximum is $7.75$, attained by $B$ and $E$ (a tie). We seed with $B$; the
final split is the same for either choice. So
$\C_{\text{new}}=\{B\}$, $\C_{\text{old}}=\{A,C,D,E\}$.

\subsection*{Reassignment round 1 (splinter $=\{B\}$)}
\begin{center}
\begin{tabular}{@{}ccccc@{}}
\toprule
$x$ & $a(x,\C_{\text{old}}\setminus\{x\})$ & $a(x,\{B\})$ & $D_x$ & \\
\midrule
$A$ & $\tfrac{3+6+11}{3}=6.6667$ & $9$  & $-2.3333$ & \\
$C$ & $\tfrac{3+9+2}{3}=4.6667$  & $7$  & $-2.3333$ & \\
$D$ & $\tfrac{6+9+8}{3}=7.6667$  & $5$  & $\mathbf{+2.6667}$ & $\leftarrow$ max $>0$\\
$E$ & $\tfrac{11+2+8}{3}=7.0000$ & $10$ & $-3.0000$ & \\
\bottomrule
\end{tabular}
\end{center}
$D_D=+2.6667>0$, so \textbf{move $D$} into the splinter:
$\C_{\text{new}}=\{B,D\}$, $\C_{\text{old}}=\{A,C,E\}$.

\subsection*{Reassignment round 2 (splinter $=\{B,D\}$)}
\begin{center}
\begin{tabular}{@{}cccc@{}}
\toprule
$x$ & $a(x,\C_{\text{old}}\setminus\{x\})$ & $a(x,\{B,D\})$ & $D_x$\\
\midrule
$A$ & $\tfrac{3+11}{2}=7.0$ & $\tfrac{9+6}{2}=7.5$  & $-0.5$\\
$C$ & $\tfrac{3+2}{2}=2.5$  & $\tfrac{7+9}{2}=8.0$  & $-5.5$\\
$E$ & $\tfrac{11+2}{2}=6.5$ & $\tfrac{10+8}{2}=9.0$ & $-2.5$\\
\bottomrule
\end{tabular}
\end{center}
Now $\max_x D_x=-0.5\le0$, so we \textbf{stop}. The root splits as
\[
\{A,B,C,D,E\}\;\longrightarrow\;\{A,C,E\}\ \mid\ \{B,D\}\qquad(\text{at diameter }11).
\]

\section{Split 2 --- Divide $\{A,C,E\}$}
The two current clusters are $\{A,C,E\}$ (diameter $\max(3,11,2)=11$) and
$\{B,D\}$ (diameter $5$). The larger diameter is $\{A,C,E\}$, so we split it.
Average dissimilarities within $\{A,C,E\}$:
\[
a(A)=\tfrac{3+11}{2}=7.0,\qquad a(C)=\tfrac{3+2}{2}=2.5,\qquad
a(E)=\tfrac{11+2}{2}=6.5 .
\]
Seed with $A$ (largest). Splinter $=\{A\}$, old $=\{C,E\}$:
\[
D_C=a(C,\{E\})-d(C,A)=2-3=-1,\qquad
D_E=a(E,\{C\})-d(E,A)=2-11=-9 .
\]
Both $\le0$, so split immediately:
\[
\{A,C,E\}\;\longrightarrow\;\{C,E\}\ \mid\ \{A\}\qquad(\text{at diameter }11).
\]

\section{Splits 3 and 4 --- The Two-Element Clusters}
Only two-object clusters remain to break; each splits trivially into its two
singletons at its own diameter:
\[
\{B,D\}\;\longrightarrow\;\{B\}\mid\{D\}\ (\text{diam }5),\qquad
\{C,E\}\;\longrightarrow\;\{C\}\mid\{E\}\ (\text{diam }2).
\]

\section{Split Summary and Divisive Tree}
\begin{center}
\begin{tabular}{@{}clcc@{}}
\toprule
Split & Cluster divided & Diameter & Resulting parts\\
\midrule
1 & $\{A,B,C,D,E\}$ & $11$ & $\{A,C,E\}$, $\{B,D\}$\\
2 & $\{A,C,E\}$     & $11$ & $\{A\}$, $\{C,E\}$\\
3 & $\{B,D\}$       & $5$  & $\{B\}$, $\{D\}$\\
4 & $\{C,E\}$       & $2$  & $\{C\}$, $\{E\}$\\
\bottomrule
\end{tabular}
\end{center}

\begin{figure}[ht]
\centering
\begin{tikzpicture}[
   every node/.style={draw,rounded corners,inner sep=3pt,font=\footnotesize},
   edge/.style={-{Latex[length=2mm]},draw,thick}]
\node (root) at (4,0)      {$\{A,B,C,D,E\}$};
\node (ace)  at (2,-1.7)   {$\{A,C,E\}$};
\node (bd)   at (6,-1.7)   {$\{B,D\}$};
\node (a)    at (0.7,-3.4) {$\{A\}$};
\node (ce)   at (3,-3.4)   {$\{C,E\}$};
\node (b)    at (5,-3.4)   {$\{B\}$};
\node (dd)   at (7,-3.4)   {$\{D\}$};
\node (c)    at (2.4,-5.1) {$\{C\}$};
\node (e)    at (3.6,-5.1) {$\{E\}$};
\draw[edge] (root)--(ace);  \draw[edge] (root)--(bd);
\draw[edge] (ace)--(a);     \draw[edge] (ace)--(ce);
\draw[edge] (bd)--(b);      \draw[edge] (bd)--(dd);
\draw[edge] (ce)--(c);      \draw[edge] (ce)--(e);
\end{tikzpicture}
\caption{Divisive (DIANA) tree, built \emph{top-down}: the root is split first
(split~1), then its parts, and so on. Arrows point from a parent cluster to the
two clusters it divides into.}
\label{fig:dianatree}
\end{figure}

\begin{remark}[DIANA vs.\ AGNES on this data]
The nested structure in Figure~\ref{fig:dianatree} is \emph{identical} to the
single-linkage agglomerative tree of Figure~\ref{fig:singdend}: both give
$\bigl(\{A\},\{C,E\}\bigr),\{B,D\}$. Here top-down and bottom-up agree, but in
general they need not: divisive methods consider the \emph{global} structure of
a cluster before cutting, whereas agglomerative methods only ever look at the
closest pair, so on other datasets the two families can produce different trees.
\end{remark}

%==============================================================================
\chapter{Choosing Clusters, Pros \& Cons, and Code}
%==============================================================================

\section{Cutting the Dendrogram}
To obtain a flat clustering, cut the dendrogram horizontally:
\begin{itemize}[nosep]
  \item Cut at a chosen height $h$: all links above $h$ are severed.
  \item Or cut to obtain a desired number $k$ of clusters.
  \item A common heuristic is to cut at the largest vertical ``gap''
        (the biggest jump between successive merge heights), since a large gap
        signals that further merging joins genuinely dissimilar groups.
\end{itemize}
For the single-linkage tree (Figure~\ref{fig:singdend}) the biggest gap is
between heights $3$ and $5$, suggesting $k=2$: $\{A,C,E\}$ and $\{B,D\}$.

\section{Advantages and Disadvantages}
\begin{center}
\begin{tabular}{@{}p{6.4cm}p{6.4cm}@{}}
\toprule
\textbf{Advantages} & \textbf{Disadvantages}\\
\midrule
No need to pre-specify $k$. & Time $\mathcal{O}(m^3)$ (naively), memory
$\mathcal{O}(m^2)$; poor for very large $m$.\\
Dendrogram gives clusters at all scales. & Merges are \emph{greedy} and
irreversible---an early bad merge cannot be undone.\\
Deterministic (no random initialization). & Sensitive to the choice of distance
and linkage.\\
Can capture non-spherical shapes (single linkage). & Single linkage suffers
from chaining; sensitive to noise/outliers.\\
\bottomrule
\end{tabular}
\end{center}

\section{Comparison with $k$-means}
\begin{center}
\begin{tabular}{@{}p{6.4cm}p{6.4cm}@{}}
\toprule
\textbf{Hierarchical} & \textbf{$k$-means}\\
\midrule
$k$ chosen \emph{after} building the tree. & $k$ fixed \emph{before} running.\\
Deterministic. & Depends on random initial centroids.\\
$\mathcal{O}(m^2)$--$\mathcal{O}(m^3)$; small/medium data. & $\mathcal{O}(kmt)$;
scales to large data.\\
Output: nested dendrogram. & Output: $k$ flat clusters.\\
\bottomrule
\end{tabular}
\end{center}

\section{Reproducing the Example in Python}
\begin{lstlisting}[style=python,caption={Agglomerative clustering with SciPy.}]
import numpy as np
from scipy.cluster.hierarchy import linkage, dendrogram, fcluster
from scipy.spatial.distance import squareform

# Symmetric distance matrix for A,B,C,D,E (same as the notes)
D = np.array([[ 0, 9, 3, 6,11],
              [ 9, 0, 7, 5,10],
              [ 3, 7, 0, 9, 2],
              [ 6, 5, 9, 0, 8],
              [11,10, 2, 8, 0]], dtype=float)

condensed = squareform(D)                 # SciPy wants the condensed form
Z = linkage(condensed, method='single')   # try 'complete', 'average', 'ward'
print(Z)                                   # merge list: [idx1, idx2, height, size]

labels = ['A','B','C','D','E']
dendrogram(Z, labels=labels)               # plot the tree
clusters = fcluster(Z, t=2, criterion='maxclust')   # cut into 2 clusters
print(dict(zip(labels, clusters)))         # -> {A,C,E} vs {B,D}
\end{lstlisting}
With \texttt{method='single'} SciPy reports merge heights $2,3,5,6$---exactly
the by-hand result of Chapter~\ref{ch:single}.

%==============================================================================
\chapter{Exercises}
%==============================================================================

\section{Solved Problems}

\begin{example}[Euclidean proximity matrix]
Compute the Euclidean distance between $P_1=(1,1)$ and $P_2=(4,5)$.\\
\textbf{Solution.} $d=\sqrt{(4-1)^2+(5-1)^2}=\sqrt{9+16}=\sqrt{25}=5.$
\end{example}

\begin{example}[First merge]
For the matrix \eqref{eq:D0}, which pair merges first and at what height under
\emph{any} linkage?\\
\textbf{Solution.} The very first merge only looks at the raw minimum entry,
which is $D_{CE}=2$; so $C,E$ merge first at height $2$ regardless of linkage.
\end{example}

\begin{example}[Single vs complete update]
Given $d(C,X)=3$ and $d(E,X)=11$, find $d((CE),X)$ under single and complete
linkage.\\
\textbf{Solution.} Single: $\min(3,11)=3$. Complete: $\max(3,11)=11$.
\end{example}

\begin{example}[Cutting for $k=3$]
Using the complete-linkage dendrogram (Figure~\ref{fig:compdend}), list the
three clusters obtained by cutting just below height $9$.\\
\textbf{Solution.} The merge of $A$ with $(BD)$ happens \emph{at} height $9$,
so a cut just below $9$ severs it along with everything above. The remaining
clusters are $\{A\}$, $\{B,D\}$ and $\{C,E\}$ --- three clusters, as required.
\end{example}

\section{Unsolved Problems}

\begin{exercise}[Build a proximity matrix]
Given points $A(1,1),B(2,1),C(4,3),D(5,4)$ in $\R^2$, compute the full
$4\times4$ Euclidean proximity matrix (round to two decimals).
\end{exercise}

\begin{exercise}[Single linkage by hand]
Run agglomerative single-linkage clustering on your matrix from Exercise~8.1,
list every merge with its height, and sketch the dendrogram.
\end{exercise}

\begin{exercise}[Complete linkage]
Repeat Exercise~8.2 with complete linkage. Do the trees differ? Explain why.
\end{exercise}

\begin{exercise}[Average linkage]
For the matrix \eqref{eq:D0}, verify by hand that the average-linkage merge
between $(ACE)$ and $(BD)$ occurs at height $49/6$.
\end{exercise}

\begin{exercise}[Lance--Williams]
Using the Lance--Williams coefficients for complete linkage
($\alpha_i=\alpha_j=\tfrac12,\ \beta=0,\ \gamma=\tfrac12$), show that
\eqref{eq:lw} reduces to $\max(d_{ik},d_{jk})$.
\end{exercise}

\begin{exercise}[Conceptual]
Explain the ``chaining'' phenomenon of single linkage and give one dataset
shape where it is an advantage and one where it is a drawback.
\end{exercise}

\begin{exercise}[Divisive]
Describe, in pseudocode, a top-down divisive scheme that at each step splits a
cluster by running $2$-means on it. What is its cost relative to agglomerative?
\end{exercise}

\begin{exercise}[Choosing $k$]
Given merge heights $2,3,5,6$ (single linkage) and $2,5,9,11$ (complete
linkage) on the same data, use the ``largest gap'' rule to recommend $k$ for
each and comment on the difference.
\end{exercise}

%==============================================================================
\appendix
\chapter{Notation Summary}
%==============================================================================
\begin{center}
\begin{tabular}{@{}ll@{}}
\toprule
Symbol & Meaning\\
\midrule
$m$ & number of objects to cluster\\
$\vx^{(i)}\in\R^n$ & the $i$-th object (feature vector)\\
$d(\va,\vb)$ & distance between two points\\
$\mathbf D$ & proximity (distance) matrix, $m\times m$\\
$\C_i$ & the $i$-th cluster; $|\C_i|$ its size\\
$d(\C_i,\C_j)$ & inter-cluster distance (linkage)\\
$\bm\mu_i$ & centroid of cluster $\C_i$\\
$\alpha_i,\alpha_j,\beta,\gamma$ & Lance--Williams coefficients\\
\bottomrule
\end{tabular}
\end{center}

\end{document}
